\documentclass[11pt]{article}
\usepackage[margin=1in]{geometry}
\usepackage{natbib}
\usepackage{hyperref}
\usepackage{booktabs}
\usepackage{enumitem}

\title{Sea-Level Rise Forecast Models Related to SPECTRUM:\\A Literature Assessment}
\author{Literature search compiled 2026-05-10; updated 2026-09-03}
\date{}

\begin{document}
\maketitle

\noindent This report identifies published frameworks and studies most
relevant to SPECTRUM, our Bayesian, observation-calibrated, component-level
GMSL forecasting framework. The original version of this report (2026-05-10)
covered only the six closest component-summation frameworks (Section~1). The
2026-09-03 update broadens the search to five additional methodological
facets of SPECTRUM: alternative component-summation frameworks (Section~2),
the semi-empirical GMST-sensitivity lineage (Section~3), deep-uncertainty and
WAIS tail-risk framing (Section~4), observational anchoring and naive
extrapolation (Section~5), and reduced-complexity discharge/calving
parameterizations analogous to our Greenland model (Section~6). Section~7
lists items published in approximately the last twelve months. A
consolidated summary and a set of specific action items for the manuscript
close the report (Section~8).

%% ========================================================================
\section{Core component-summation frameworks (2026-05-10 search)}

\subsection*{1.1 Mengel et al.\ (2016)}

\noindent\textbf{Citation:} Mengel, M., Levermann, A., Frieler, K.,
Robinson, A., Marzeion, B., \& Winkelmann, R.\ (2016). Future sea level
rise constrained by observations and long-term commitment.
\textit{Proc.\ Natl.\ Acad.\ Sci.}, 113(10), 2597--2602.\\
\textbf{DOI:} \texttt{10.1073/pnas.1500515113}

\noindent\textbf{Approach.} Each component (thermal expansion, glaciers,
Greenland, Antarctica) is represented by a simple physically motivated
response function with a small number of parameters, independently
calibrated to 20th-century observations via Bayesian inference. Projections
are generated by sampling posteriors and forcing with RCP temperature
scenarios, then summing components. Inter-component correlations arise only
through shared temperature forcing.

\noindent\textbf{Relevance: Very high.} This is the closest published
methodological parallel: per-component semi-empirical models, each
observation-calibrated, summed probabilistically. The main differences are
(i) our use of rate-space and level-space Bayesian models rather than
equilibrium response functions, (ii) our explicit deep-uncertainty
treatment of WAIS, and (iii) our use of a blending step that anchors
projections to satellite-era observed rates.

\subsection*{1.2 Kopp et al.\ (2014)}

\noindent\textbf{Citation:} Kopp, R.\,E., Horton, R.\,M., Little, C.\,M.,
Mitrovica, J.\,X., Oppenheimer, M., Rasmussen, D.\,J., Strauss, B.\,H.,
\& Tebaldi, C.\ (2014). Probabilistic 21st and 22nd century sea-level
projections at a global network of tide-gauge sites.
\textit{Earth's Future}, 2(8), 383--406.\\
\textbf{DOI:} \texttt{10.1002/2014EF000239}

\noindent\textbf{Approach.} The foundational probabilistic
component-summation framework. Individual components (thermal expansion,
glaciers, GrIS, AIS, land water storage) are each represented as
probability distributions drawn from IPCC AR5 expert assessment and process
models. Monte Carlo sampling sums components at each location, applying
gravitational fingerprints for regional projections.

\noindent\textbf{Relevance: High.} Established the ``plug-and-play''
component-summation architecture that all subsequent probabilistic
frameworks build on, including IPCC AR6. Our work shares this architecture
but replaces assessment-derived component PDFs with observation-calibrated
Bayesian posteriors, and does not address regional patterns.

\subsection*{1.3 Wong et al.\ (2017) --- BRICK}

\noindent\textbf{Citation:} Wong, T.\,E., Bakker, A.\,M.\,R.,
Ruckert, K., Applegate, P., Slangen, A.\,B.\,A., \& Keller, K.\ (2017).
BRICK v0.2, a simple, accessible, and transparent model framework for
climate and regional sea-level projections.
\textit{Geosci.\ Model Dev.}, 10, 2741--2760.\\
\textbf{DOI:} \texttt{10.5194/gmd-10-2741-2017}

\noindent\textbf{Approach.} Couples simple, physically motivated models for
each component (DOECLIM for thermal expansion, Marzeion-type scaling for
glaciers, SMB/dynamics models for GrIS and AIS) within a shared climate
state. The full system is jointly calibrated to observations using MCMC,
producing posteriors that naturally capture inter-component correlations.

\noindent\textbf{Relevance: High.} BRICK is the only published framework
that performs joint Bayesian calibration of all components simultaneously,
making it the closest analogue to our approach in terms of statistical
methodology. A key distinction is that our framework calibrates each
component independently and handles WAIS through a structured scenario
mixture rather than a process model. See also Darnell et al.\ (2025,
Section~4.2), a BRICK-based follow-on directly relevant to our
emissions-vs-tail-risk framing.

\subsection*{1.4 Kopp et al.\ (2023) --- FACTS}

\noindent\textbf{Citation:} Kopp, R.\,E., Garner, G.\,G.,
Hermans, T.\,H.\,J., et al.\ (2023). The Framework for Assessing Changes
To Sea-level (FACTS) v1.0: a platform for characterizing parametric and
structural uncertainty in future global, relative, and extreme sea-level
change. \textit{Geosci.\ Model Dev.}, 16, 7461--7489.\\
\textbf{DOI:} \texttt{10.5194/gmd-16-7461-2023}

\noindent\textbf{Approach.} The modular software platform underlying IPCC
AR6 sea level projections. Each component is a pluggable module that can use
different process models or emulators. The framework propagates both
parametric uncertainty (within a model) and structural uncertainty (across
alternative models for the same component).

\noindent\textbf{Relevance: High.} FACTS is the institutional benchmark
against which any new component-level framework will be compared. Our work
can be understood as a complementary approach that replaces the
process-model modules in FACTS with data-driven Bayesian models, trading
physical detail for tighter observational constraint and transparency about
what the data alone can support.

\subsection*{1.5 Fox-Kemper et al.\ (2021) --- IPCC AR6 Chapter 9}

\noindent\textbf{Citation:} Fox-Kemper, B., Hewitt, H.\,T., Xiao, C.,
et al.\ (2021). Ocean, Cryosphere and Sea Level Change. In
\textit{Climate Change 2021: The Physical Science Basis} (pp.\ 1211--1362).
Cambridge University Press.\\
\textbf{DOI:} \texttt{10.1017/9781009157896.011}

\noindent\textbf{Approach.} The authoritative component-level assessment,
summing thermal expansion (CMIP6), glaciers (GlacierMIP), Greenland
(ISMIP6 + emulators), Antarctica (ISMIP6 + emulators + structured expert
judgment), and land water storage. Provides both medium-confidence and
low-confidence projections across all SSP scenarios.

\noindent\textbf{Relevance: High.} This is the primary comparison target
for any new SLR projection. Our framework produces systematically higher
medians and wider uncertainty ranges than IPCC AR6 medium confidence,
largely because the IPCC ice sheet component relies on ISMIP6 process
models that may understate dynamic mass loss. The IPCC low-confidence
projections, which incorporate expert judgment, are more consistent with our
results. Törnqvist et al.\ (2025, Section~5) is a useful counterpoint here:
it shows the 1995/96 IPCC SAR mid-range projection verified to within
$\sim$1~cm over the following 30 years, so any claim that ``IPCC projections
diverge from observations'' should specify AR6 specifically rather than
IPCC assessments in general.

\subsection*{1.6 Nauels et al.\ (2017) --- MAGICC Sea Level Model}

\noindent\textbf{Citation:} Nauels, A., Meinshausen, M., Mengel, M.,
Lorbacher, K., \& Wigley, T.\,M.\,L.\ (2017). Synthesizing long-term sea
level rise projections --- the MAGICC sea level model v2.0.
\textit{Geosci.\ Model Dev.}, 10, 2495--2524.\\
\textbf{DOI:} \texttt{10.5194/gmd-10-2495-2017}

\noindent\textbf{Approach.} Emulates process-based responses for seven
sub-components (thermal expansion, glaciers, GrIS SMB, GrIS dynamics, WAIS,
EAIS, land water storage) using reduced-complexity parameterizations
calibrated to CMIP5/AR5 output. Embedded within the MAGICC climate model,
ensuring a self-consistent climate trajectory across components. Extends
projections to 2300.

\noindent\textbf{Relevance: Moderate--high.} Shares the component
granularity of our framework (seven sub-components) and the goal of
efficient probabilistic ensemble generation. However, MAGICC calibrates to
process-model output rather than directly to observations, and does not
incorporate deep uncertainty in ice sheet dynamics. Our framework's direct
observational calibration and explicit WAIS scenario structure address both
of these limitations. See also Nauels et al.\ (2025, Section~7), a recent
multi-century commitment framework from an overlapping author group.

%% ========================================================================
\section{Additional component-summation and probabilistic frameworks}

\subsection*{2.1 Le Bars \& Drijfhout (2025)}

\noindent\textbf{Citation:} Le Bars, D. \& Drijfhout, S.\ (2025). Connecting
the Past and the Future to Build Regional Sea-Level Projections.
\textit{Earth's Future}, 13, e2024EF005623.

\noindent\textbf{Approach.} The KNMI (Netherlands) national framework builds
a GMSL/regional budget (Antarctica, Greenland, glaciers, thermosteric, land
water, vertical land motion) that is deliberately anchored to start in 1995
so that it overlaps and is evaluated against the observational record before
extending to 2100.

\noindent\textbf{Relevance: High.} Of everything found in this search, this
is structurally the closest published analogue to SPECTRUM's practice of
calibrating each component against its own observational record and
checking budget closure against satellite altimetry. Worth a direct
comparison in the framework section.

\subsection*{2.2 Felikson et al.\ (2025)}

\noindent\textbf{Citation:} Felikson, D., Rounce, D.\,R., Fasullo, J.,
Rodriguez, A., Adhikari, S., Buzzanga, B., Dangendorf, S., Kopp, R.\,E.,
Lammers, R.\,B., Reager, J.\,T., Brinkerhoff, D., Csatho, B., Girotto, M.,
Hamlington, B., Ivins, E.\,R., Kumar, P., Larour, E., Nerem, R.\,S.,
Nowicki, S., Schlegel, N.-J., Tesdal, J.-E., \& Weathers, M.\ (2025).
Progress and future directions in constraining uncertainties in
sea-level projections using observations.
\textit{Nature Climate Change}, 15(10), 1039--1051.\\
\textbf{DOI:} \texttt{10.1038/s41558-025-02437-4}

\noindent\textbf{Approach.} A 22-author community perspective (co-authored
by several of the researchers behind the leading process-based and
observational SLR groups: Kopp, Nerem, Hamlington, Dangendorf, Nowicki,
among others) reviewing the main sources of uncertainty in multi-decadal to
centennial SLR projections, and arguing explicitly for using observations to
narrow them. It calls for improved emulation techniques, more systematic
quantification of uncertainty \emph{structure} in both observations and
models, longer observational records for the underlying physical processes,
and closer collaboration across physical and social sciences, framed
explicitly against the next IPCC assessment cycle.

\noindent\textbf{Relevance: Very high.} This is the single closest published
statement of SPECTRUM's underlying philosophy --- constrain component
projections with observations rather than relying on process models alone
--- from a broad, senior, cross-institutional author group rather than a
single competing framework. It is not itself a competing forecasting
method, so it does not need to be distinguished from SPECTRUM the way the
Section~1 frameworks do; instead it is best cited in the introduction or
discussion as the field-level motivation SPECTRUM directly answers. See
Section~9 (``near-perfect conceptual match'') for a fuller explanation.

\subsection*{2.3 Slangen et al.\ (2022)}

\noindent\textbf{Citation:} Slangen, A.\,B.\,A., Palmer, M.\,D.,
Camargo, C.\,M.\,L., Church, J.\,A., Edwards, T.\,L., Hermans, T.\,H.\,J.,
Hewitt, H.\,T., Garner, G.\,G., et al.\ (2022). The evolution of 21st
century sea-level projections from IPCC AR5 to AR6 and beyond.
\textit{Cambridge Prisms: Coastal Futures}, 1, e8.\\
\textbf{DOI:} \texttt{10.1017/cft.2022.8}

\noindent\textbf{Approach.} Traces how successive IPCC component-summation
frameworks (AR5, SROCC, AR6) changed their component treatments, uncertainty
structure, and headline numbers across assessment cycles.

\noindent\textbf{Relevance: Moderate--high.} Useful for positioning SPECTRUM
against the IPCC lineage as a whole rather than against AR6 alone, and for
anticipating how the next assessment cycle (which Felikson et al.\ 2025
explicitly addresses) is likely to evolve.

\subsection*{2.4 Related items noted but not detailed}

Garner et al.\ (2018, \textit{Earth's Future}, DOI
\texttt{10.1029/2018EF000991}) is an earlier meta-comparison of the
evolution of SLR projections (1980s--2018), complementary to Slangen et al.\
(2022). Robel, Seroussi \& Roe (2019, \textit{PNAS}, DOI
\texttt{10.1073/pnas.1904822116}) shows marine ice sheet instability (MISI)
mathematically skews the SLR probability distribution toward high outcomes;
it is discussed at length in Section~4.3 but is equally relevant here as the
theoretical justification for why a component-summation posterior for WAIS
should be right-skewed rather than Gaussian. Le~Cozannet, Manceau \& Rohmer
(2017, \textit{Environ.\ Res.\ Lett.}) apply possibility theory rather than
Bayesian Monte Carlo to bound probabilistic SLR projections --- a genuinely
different uncertainty formalism worth a citation if the discussion engages
with alternatives to Bayesian summation.

%% ========================================================================
\section{Semi-empirical GMST-sensitivity models}

\subsection*{3.1 Rahmstorf (2007) and the rate-vs-level question}

\noindent\textbf{Citation:} Rahmstorf, S.\ (2007). A Semi-Empirical
Approach to Projecting Future Sea-Level Rise. \textit{Science}, 315,
368--370.\\
\textbf{DOI:} \texttt{10.1126/science.1135456} --- already cited in the
manuscript as \texttt{Rahmstorf2007}.

\noindent\textbf{Approach.} Regresses the rate of GMSLR linearly on the
\emph{level} of GMST, $dS/dt = a\,T$, calibrated on the 20th-century
instrumental record; the founding paper of the semi-empirical class the
manuscript already discusses.

\noindent\textbf{Citation:} Vermeer, M. \& Rahmstorf, S.\ (2009). Global
sea level linked to global temperature. \textit{PNAS}, 106, 21527--21532.\\
\textbf{DOI:} \texttt{10.1073/pnas.0907765106} --- already present in
\texttt{references.bib} as \texttt{Vermeer2009} but \textbf{not currently
cited in the manuscript text}.

\noindent\textbf{Approach.} Extends Rahmstorf (2007) by adding a rate-of-warming
term, $dS/dt = a\,T + b\,dT/dt$, so that GMSLR responds to both the level and
the rate of change of temperature.

\noindent\textbf{Relevance: Very high.} This is the direct published
precedent for the open rate-vs-level question already flagged in project
notes (does SLR respond to $T$ or also $dT/dt$?). It should be cited
alongside Rahmstorf (2007) wherever the manuscript discusses the
semi-empirical class, and is the natural anchor for any explicit discussion
of that open question.

\subsection*{3.2 Grinsted \& Christensen (2021) --- transient sensitivity}

\noindent\textbf{Citation:} Grinsted, A. \& Christensen, J.\,H.\ (2021). The
transient sensitivity of sea level rise. \textit{Ocean Science}, 17,
181--186.\\
\textbf{DOI:} \texttt{10.5194/os-17-181-2021} --- already in the local
Zotero library, not currently cited.

\noindent\textbf{Approach.} Reframes the rate-vs-level question again: shows
that on century timescales, GMSLR rate correlates better with
\emph{accumulated (cumulative) warming} than with instantaneous GMST,
defining a ``transient sea-level sensitivity'' (TSLS) distinct from the
millennial-scale equilibrium sensitivity. The observed TSLS
($0.40 \pm 0.05$~m~$^\circ$C$^{-1}$~century$^{-1}$) exceeds CMIP-based
estimates (0.27--0.39). A companion paper, Grinsted et al.\ (2022,
\textit{Earth's Future}, DOI \texttt{10.1029/2022EF002696}, in Zotero),
extends the same diagnostic to CMIP6 GCMs directly.

\noindent\textbf{Relevance: High.} A third, genuinely distinct framing of
the level-vs-rate question (cumulative-forcing response) that the manuscript
should engage with directly if it takes a position on this question,
alongside Rahmstorf (2007) and Vermeer \& Rahmstorf (2009).

\subsection*{3.3 Critiques of the semi-empirical approach}

The manuscript currently attributes IPCC AR5's ``low confidence''
characterization of semi-empirical models solely to Church et al.\ (2013,
already cited as \texttt{Church2013}), but that assessment synthesized
specific empirical critiques not currently cited.
The most relevant, none yet in the local Zotero library:

\begin{itemize}[leftmargin=*]
\item Holgate, S., Jevrejeva, S., Woodworth, P., \& Brewer, S.\ (2007).
Comment on ``A Semi-Empirical Approach to Projecting Future Sea-Level
Rise.'' \textit{Science}, 317, 1866 (DOI \texttt{10.1126/science.1140942}),
with Rahmstorf's reply in the same issue (DOI
\texttt{10.1126/science.1141283}). Argues no robust linear GMST-rate
relationship is detectable in the data used; Rahmstorf's reply defends
split-sample predictive skill.
\item Orlić, M. \& Pasarić, Z.\ (2013). Semi-empirical versus
process-based sea-level projections for the twenty-first century.
\textit{Nat.\ Clim.\ Change}, 3, 735--738 (DOI \texttt{10.1038/nclimate1877}),
and (2015), Some Pitfalls of the Semiempirical Method Used to Project Sea
Level Rise. \textit{J.\ Climate}, 28, 3779--3785 (DOI
\texttt{10.1175/JCLI-D-14-00696.1}). Shows semi-empirical projections are
highly sensitive to which dynamical terms and detrending/response-time
choices are included; a ``corrected'' estimate converges toward
process-based values.
\item Schmith, T., Johansen, S., \& Thejll, P.\ (2012). Statistical
Analysis of Global Surface Temperature and Sea Level Using Cointegration
Methods. \textit{J.\ Climate}, 25, 7822--7833. Finds temperature and sea
level share a common stochastic trend, but the 1940s warming episode
deviates from the fitted relationship, attributed to internal ocean
variability rather than external forcing --- a direct challenge to
stationarity assumptions underlying extrapolation.
\item Grassi, S., Hillebrand, E., \& Ventosa-Santaularia, D.\ (2013). The
statistical relation of sea-level and temperature revisited. \textit{Dyn.\
Atmos.\ Oceans} (special issue). Re-estimates Rahmstorf's proportionality
coefficient with methods appropriate for two co-trending series; the
sensitivity remains at only marginal statistical significance and the 90\%
CI on 2100 GMSLR widens to 15--150~cm.
\end{itemize}

\noindent\textbf{Relevance: Moderate.} Not competing frameworks for
SPECTRUM, but the manuscript's brief mention of semi-empirical models should
cite 1--2 of these directly rather than relying on the AR5 assessment alone
as a secondhand summary of the critique.

%% ========================================================================
\section{Deep uncertainty, MISI/MICI, and WAIS tail-risk framing}

\subsection*{4.1 Naughten, Holland \& De Rydt (2023)}

\noindent\textbf{Citation:} Naughten, K.\,A., Holland, P.\,R., \&
De Rydt, J.\ (2023). Unavoidable future increase in West Antarctic
ice-shelf melting over the twenty-first century. \textit{Nature Climate
Change}, 13, 1222--1228.\\
\textbf{DOI:} \texttt{10.1038/s41558-023-01818-x} --- already in the local
Zotero library.

\noindent\textbf{Approach.} A regional ocean model of the Amundsen Sea
embayment shows ice-shelf melt accelerating at similar rates across all
emissions scenarios tested this century, i.e., that ocean-driven WAIS melt
through roughly 2100 is largely already locked in.

\noindent\textbf{Relevance: Very high.} This is the strongest existing,
independent, dynamical-modeling support for SPECTRUM's headline claim that
the high-risk tail is largely insensitive to emissions scenario. It should
be a central citation for that claim.

\subsection*{4.2 Darnell et al.\ (2025) --- a counterpoint to engage directly}

\noindent\textbf{Citation:} Darnell, C., Rennels, L., Errickson, F.,
Wong, T.\,E., \& Srikrishnan, V.\ (2025). The interplay of future emissions
and geophysical uncertainties for projections of sea-level rise.
\textit{Nature Climate Change}.\\
\textbf{DOI:} \texttt{10.1038/s41558-025-02457-0}

\noindent\textbf{Approach.} Using the calibrated BRICK model (Wong et al.\
2017, Section~1.3), decomposes SLR projection variance into emissions-driven
and geophysical (largely ice-sheet) components as a function of forecast
horizon. Finds emissions-scenario uncertainty dominates in the near-to-mid
term (specifically the 2065--2075 window), with geophysical/tipping-point
uncertainty in Antarctica becoming dominant only later.

\noindent\textbf{Relevance: Very high --- requires a direct response.} This
paper's headline decomposition is close to a rival claim to SPECTRUM's: it
argues emissions choice \emph{does} materially affect the SLR distribution
over multi-decadal horizons before geophysical uncertainty takes over,
whereas SPECTRUM emphasizes that emissions cuts move the median but leave
the tail (which we attribute predominantly to WAIS) largely unchanged. This
surfaced independently from two separate literature searches run for this
report, which is a reasonable signal that it is load-bearing rather than a
minor citation. The manuscript should engage with it explicitly --- for
example, by clarifying that our claim is about the shape/tail of the
distribution at a fixed horizon (end of century) rather than about
total-variance decomposition across horizons, which are not necessarily in
conflict, but the distinction should be stated rather than left for a
referee to raise.

\subsection*{4.3 MISI/MICI physics and structural uncertainty}

\begin{itemize}[leftmargin=*]
\item Robel, A.\,A., Seroussi, H., \& Roe, G.\,H.\ (2019). Marine ice sheet
instability amplifies and skews uncertainty in projections of future
sea-level rise. \textit{PNAS}, 116(30), 14887--14892 (DOI
\texttt{10.1073/pnas.1904822116}; in Zotero). Stochastic perturbation
analysis showing MISI dynamics structurally produce a right-skewed SLR
distribution --- theoretical grounding for treating WAIS with a fat right
tail rather than choosing one ad hoc.
\item DeConto, R.\,M. \& Pollard, D.\ (2016). Contribution of Antarctica to
past and future sea-level rise. \textit{Nature}, 531, 591--597 (DOI
\texttt{10.1038/nature17145}; in Zotero). The foundational MICI/hydrofracturing
model producing multi-meter 21st-century Antarctic contributions under high
emissions.
\item Edwards, T.\,L., et al.\ (2019). Revisiting Antarctic ice loss due to
marine ice-cliff instability. \textit{Nature}, 566, 58--64 (DOI
\texttt{10.1038/s41586-019-0901-4}). A statistical emulator with properly
propagated parametric uncertainty concludes the original DeConto \& Pollard
projections likely overestimated central/high-end SLR.
\item Morlighem, M., Goldberg, D., Barnes, J.\,M., et al.\ (2024). The West
Antarctic Ice Sheet may not be vulnerable to marine ice cliff instability
during the 21st century. \textit{Science Advances}, 10, eado7794 (in
Zotero). A multi-model intercomparison arguing MICI is unlikely to activate
before 2100 under realistic near-term forcing.
\item Bassis, J.\,N., Berg, B., Crawford, A.\,J., \& Benn, D.\,I.\ (2021).
Transition to marine ice cliff instability controlled by ice thickness
gradients and velocity. \textit{Science}, 372, 1342--1344. Refines the
physical trigger condition for MICI, showing the transition to runaway
calving is more physically constrained than the original parameterization.
\end{itemize}

\noindent\textbf{Relevance: High, but mixed.} Together these papers show
that whether and when MICI activates this century is itself an actively
contested question in the ice-sheet-dynamics literature. Edwards (2019) and
Morlighem (2024) are important tempering counter-evidence that should be
cited alongside DeConto \& Pollard when justifying the probability weight
assigned to SPECTRUM's fast/MICI-type WAIS branch, rather than treating MICI
risk as settled.

\subsection*{4.4 Deep-uncertainty decision framing and near-term predictability}

\begin{itemize}[leftmargin=*]
\item Bakker, A.\,M.\,R., Wong, T.\,E., Ruckert, K.\,L., \& Keller, K.\
(2017). Sea-level projections representing the deeply uncertain
contribution of the West Antarctic ice sheet. \textit{Sci.\ Rep.}, 7, 3880
(DOI \texttt{10.1038/s41598-017-04134-5}). Explicitly frames WAIS via
deep-uncertainty decision analysis rather than a process model --- the
closest existing methodological precedent for SPECTRUM's WAIS scenario
mixture and terminology.
\item Kopp, R.\,E., DeConto, R.\,M., Bader, D.\,A., et al.\ (2017). Evolving
Understanding of Antarctic Ice-Sheet Physics and Ambiguity in Probabilistic
Sea-Level Projections. \textit{Earth's Future}, 5, 1217--1233 (in Zotero).
Shows that alternative structural assumptions about Antarctic physics (with
vs.\ without MICI-type processes) move probabilistic SLR distributions far
more than emissions-scenario choice does --- close precedent for
SPECTRUM's ``orthogonal dimension of risk'' framing.
\item McCormack, F.\,S., et al.\ (2026). Emergent decadal predictability in
Antarctic contribution to sea level rise. \textit{Nature}, 654, 609--613
(DOI \texttt{10.1038/s41586-026-10614-4}). Finds the 2025 Antarctic
mass-loss rate is strongly predictive of the rate for the following 3--5
decades across models and scenarios, with predictability breaking down only
late-century as marine ice-sheet feedbacks emerge --- independent
dynamical-modeling evidence converging with SPECTRUM's rate-anchoring
strategy and its finding that most forecast uncertainty is late-century and
WAIS-driven.
\end{itemize}

%% ========================================================================
\section{Observational anchoring, nowcasting, and naive extrapolation}

\subsection*{5.1 Closest precedents to the naive-extrapolation baseline}

\noindent\textbf{Citation:} Nerem, R.\,S., Frederikse, T., \&
Hamlington, B.\,D.\ (2022). Extrapolating Empirical Models of
Satellite-Observed Global Mean Sea Level to Estimate Future Sea Level
Change. \textit{Earth's Future}, 10, e2021EF002290 --- already in the local
Zotero library.

\noindent\textbf{Approach.} Directly extrapolates a linear-plus-quadratic
fit of satellite altimetry GMSL to 2050 and compares it against IPCC AR6,
finding rough consistency but flagging that the extrapolation assumes
constant rate and acceleration.

\noindent\textbf{Relevance: Very high.} The closest existing precedent to
SPECTRUM's own naive-extrapolation baseline; should be cited and explicitly
differentiated (SPECTRUM extends the horizon to 2100 and treats the naive
forecast as one rung in an explicit model-complexity hierarchy rather than
a standalone estimate).

\noindent Also relevant: Nerem et al.\ (2018, \textit{PNAS}, DOI
\texttt{10.1073/pnas.1717312115}), the original detection/quantification of
GMSLR acceleration from the altimetry record ($\sim$0.084~mm~yr$^{-2}$);
Hamlington et al.\ (2024, \textit{Comm.\ Earth \& Environ.}, DOI
\texttt{10.1038/s43247-024-01761-5}), which quantifies the doubling of GMSLR
rate over 1993--2023 underlying the manuscript's own doubling-time
discussion; and Dangendorf et al.\ (2019, \textit{Nat.\ Clim.\ Change}),
which shows the acceleration signal predates the satellite era, supporting
the robustness of extrapolating it.

\subsection*{5.2 Observational constraint as a general method}

\noindent\textbf{Citation:} Ribes, A., Qasmi, S., \& Gillett, N.\,P.\
(2021). Making climate projections conditional on historical observations.
\textit{Science Advances}, 7, eabc0671.\\
\textbf{DOI:} \texttt{10.1126/sciadv.abc0671}

\noindent\textbf{Approach.} A general statistical framework, outside the SLR
literature, for constraining long-term CMIP-style warming projections using
the full historical observational record, narrowing projection uncertainty
by roughly a factor of 2--3.

\noindent\textbf{Relevance: High.} The clearest methodological ancestor,
outside sea level specifically, for SPECTRUM's rate-space blending step.
Useful as a citation establishing that observationally-constrained
projection is an established technique in climate science more broadly, not
an ad hoc device specific to this study.

\subsection*{5.3 A necessary nuance}

\noindent\textbf{Citation:} Törnqvist, T.\,E., Conrad, C.\,P.,
Dangendorf, S., \& Hamlington, B.\,D.\ (2025). Evaluating IPCC Projections
of Global Sea-Level Change From the Pre-Satellite Era. \textit{Earth's
Future}, 13, e2025EF006533.

\noindent\textbf{Approach.} Finds that the IPCC Second Assessment Report
(1995/96) mid-range GMSL projection verified to within $\sim$1~cm accuracy
over the following 30 years.

\noindent\textbf{Relevance: High --- a needed qualifier.} SPECTRUM's
critique of process-model divergence from observations is specifically a
critique of \emph{AR6}, not of IPCC assessments in general; this paper
shows an earlier IPCC vintage performed very well. The manuscript's
introduction should make this scope explicit rather than let the broader
claim stand unqualified (cross-referenced in Section~1.5).

\subsection*{5.4 Simplicity-in-forecasting literature}

Beyond Green \& Armstrong (2015, already cited as \texttt{Green2015}),
Makridakis, S., Spiliotis, E., \& Assimakopoulos, V.\ (2020), The M4
Competition: 100,000 time series and 61 forecasting methods,
\textit{Int.\ J.\ Forecasting}, 36, 54--74 (DOI
\texttt{10.1016/j.ijforecast.2019.04.014}), is a large-scale empirical
forecasting competition confirming that simple statistical and hybrid
methods remain highly competitive against complex ones --- a useful second,
independent citation for the general model-parsimony argument, alongside
Held (2005, Section~6.4).

%% ========================================================================
\section{Reduced-complexity discharge and calving parameterizations}

\subsection*{6.1 Closest precedent to the Greenland delayed-response model}

\noindent\textbf{Citation:} Levermann, A., Winkelmann, R., Nowicki, S.,
et al.\ (2014). Projecting Antarctic ice discharge using response functions
from SeaRISE ice-sheet models. \textit{Earth Syst.\ Dynam.}, 5(2), 271--293.\\
\textbf{DOI:} \texttt{10.5194/esd-5-271-2014}

\noindent\textbf{Approach.} Derives linear, time-lagged response
(convolution) functions relating ice discharge to ocean-driven basal-melt
forcing, calibrated against full ice-sheet-model experiments rather than
resolved at run time; this discharge parameterization underlies LARMIP and
subsequent ISMIP6-adjacent emulators.

\noindent\textbf{Relevance: High.} The closest published precedent to
SPECTRUM's own linear, time-lagged discharge model,
$D = \gamma\,T_{\mathrm{ocean}}(t-\delta) + r_0$
(Eq.~\ref{eq:greenland_rate} in the main text), though Levermann et al.\
calibrate against model output rather than direct observations as SPECTRUM
does.

\subsection*{6.2 Plume-theory melt-rate scaling underlying the discharge sensitivity}

\noindent\textbf{Citation:} Slater, D.\,A., Goldberg, D.\,N.,
Nienow, P.\,W., \& Cowton, T.\,R.\ (2016). Scalings for Submarine Melting at
Tidewater Glaciers from Buoyant Plume Theory. \textit{J.\ Phys.\
Oceanogr.}, 46(6), 1839--1855.

\noindent\textbf{Approach.} Derives the canonical plume-theory melt-rate
scaling relating subglacial discharge and ocean thermal forcing to submarine
melt rate.

\noindent\textbf{Relevance: Moderate--high.} The theoretical basis for the
``Slater/Wood-type'' melt sensitivities already cited in the manuscript's
dimensional-scaling argument (Eqs.~\ref{eq:melt_rate}--\ref{eq:gamma_aggregate});
several downstream applications of this scaling (Xu et al.\ 2013,
Rignot et al.\ 2016, Slater et al.\ 2019) are already cited in the
manuscript and in Zotero.

\subsection*{6.3 A candidate primary citation for the model-complexity argument}

\noindent\textbf{Citation:} Aschwanden, A., Bartholomaus, T.\,C.,
Brinkerhoff, D.\,J., \& Truffer, M.\ (2021). Brief communication: A roadmap
towards credible projections of ice sheet contribution to sea level.
\textit{The Cryosphere}, 15, 5705--5715. \\
\textbf{DOI:} \texttt{10.5194/tc-15-5705-2021} --- already cited in the
manuscript as \texttt{Aschwanden2021}, but currently only for the specific
claim that process models diverge from hindcasts; it is also a more general
statement of the manuscript's model-complexity argument.

\noindent\textbf{Approach.} Argues that current ice-sheet models are
systematically under-constrained relative to their degrees of freedom, and
that projection credibility requires matching model complexity to the
calibration data actually available.

\noindent\textbf{Relevance: High.} Beyond its existing citation, this paper
is close to a direct statement of SPECTRUM's central methodological
argument and could be cited a second time, alongside Green (2015) and Held
(2005), specifically where the manuscript makes the general
complexity-vs.-data-constraint case.

\subsection*{6.4 General model-hierarchy literature}

\noindent\textbf{Citation:} Held, I.\,M.\ (2005). The Gap Between Simulation
and Understanding in Climate Modeling. \textit{Bull.\ Amer.\ Meteor.\
Soc.}, 86(11), 1609--1614.

\noindent\textbf{Approach.} The canonical argument for a ``hierarchy of
models'' bridging idealized, well-understood low-order models and
comprehensive simulation models, arguing that scientific understanding (and
often forecast robustness) is better served by simpler models than by
ever-more-complex ones.

\noindent\textbf{Relevance: High.} The closest general antecedent, outside
the SLR or ice-sheet literature specifically, to the manuscript's own
``hierarchy of model complexity'' framing (naive $\rightarrow$ semi-empirical
$\rightarrow$ component-level). Pairs naturally with Green (2015) as a
second anchor citation for that argument; see also Jeevanjee, N.,
Hassanzadeh, P., Hill, S., \& Sheshadri, A.\ (2017), A perspective on
climate model hierarchies, \textit{JAMES}, 9(4), 1760--1771 (DOI
\texttt{10.1002/2017MS001038}), a more formally developed follow-on.

%% ========================================================================
\section{Items published in approximately the last twelve months (Sept.\ 2025--Sept.\ 2026)}

Most items in this window are cross-referenced above by topic (Felikson et
al.\ 2025, Darnell et al.\ 2025, Naughten-adjacent items, Törnqvist et al.\
2025, McCormack et al.\ 2026). Additional recent items not detailed
elsewhere:

\begin{itemize}[leftmargin=*]
\item Perrette, M. \& Mengel, M.\ (2025). Relative sea level projections
constrained by historical trends at tide gauge sites. \textit{Science
Advances} (DOI \texttt{10.1126/sciadv.ado4506}). Generates local SLR
projections directly from tide-gauge/GPS/altimetry historical trends rather
than from component decomposition; global medians/ranges broadly consistent
with AR6. \textbf{Relevance: High} as a genuinely different observation-driven
alternative to component summation, worth a direct contrast in the framework
discussion.
\item Coulon, V., Klose, A.\,K., Edwards, T., Turner, F., Pattyn, F., \&
Winkelmann, R.\ (2025). From short-term uncertainties to long-term
certainties in the future evolution of the Antarctic Ice Sheet.
\textit{Nature Communications}, 16, 10385 (DOI
\texttt{10.1038/s41467-025-66178-w}). Argues short-term (decadal) and
long-term (centennial) Antarctic ice-sheet uncertainty are largely
decoupled --- converges with McCormack et al.\ (2026, Section~4.4) and
supports SPECTRUM's near-term rate-anchoring logic.
\item Winkelmann, R., Garbe, J., Donges, J.\,F., Albrecht, T., et al.\
(2026). Mapping tipping risks from Antarctic ice basins under global
warming. \textit{Nature Climate Change} (DOI
\texttt{10.1038/s41558-025-02554-0}). Basin-by-basin analysis of gradual
vs.\ threshold/tipping-point ice loss across Antarctica under different
warming levels; complements SPECTRUM's WAIS scenario mixture with a
basin-resolved tipping-point framing.
\item Nauels, A., Nicholls, Z., Möller, T., Hermans, T.\,H.\,J.,
Mengel, M., Kloenne, U., Smith, C., Slangen, A.\,B.\,A., \&
Palmer, M.\,D.\ (2025). Multi-century global and regional sea-level rise
commitments from cumulative greenhouse gas emissions in the coming decades.
\textit{Nature Climate Change}, 15(11) (DOI
\texttt{10.1038/s41558-025-02452-5}) --- already in Zotero. A direct
descendant of Nauels et al.\ (2017, Section~1.6); quantifies how near-term
emissions choices lock in 0.3--0.8~m of GMSLR commitment by 2300, a useful
comparator for SPECTRUM's own emissions-vs-tail-risk framing.
\end{itemize}

\noindent No new IPCC-adjacent government or task-force assessment (NOAA
Interagency Task Force, UK Climate Change Committee, Dutch Delta Programme)
was identified in this twelve-month window.

%% ========================================================================
\section{Summary}

\begin{table}[h]
\centering\small
\begin{tabular}{@{}llll@{}}
\toprule
\textbf{Study} & \textbf{Calibration source} & \textbf{Components} & \textbf{Relevance} \\
\midrule
\multicolumn{4}{@{}l}{\textit{Core component-summation frameworks (Sec.~1)}} \\
Mengel et al.\ (2016)      & Observations (indep.)   & 4    & Very high \\
Kopp et al.\ (2014)        & IPCC assessment         & 5    & High \\
Wong et al.\ (2017)        & Observations (joint)    & 5    & High \\
Kopp et al.\ (2023)        & Process models          & 5+   & High \\
Fox-Kemper et al.\ (2021)  & CMIP6 + ISMIP6          & 5    & High \\
Nauels et al.\ (2017)      & CMIP5 emulation         & 7    & Moderate--high \\
\multicolumn{4}{@{}l}{\textit{Additional frameworks (Sec.~2)}} \\
Le Bars \& Drijfhout (2025)& Observations (indep.)   & 6    & High \\
Felikson et al.\ (2025)    & --- (perspective)       & n/a  & Very high \\
\multicolumn{4}{@{}l}{\textit{Semi-empirical / non-component (Sec.~3, 5)}} \\
Rahmstorf (2007)           & Observations (GMST)     & 1    & Very high \\
Vermeer \& Rahmstorf (2009)& Observations (GMST)     & 1    & Very high \\
Grinsted \& Christensen (2021) & Observations (GMST) & 1    & High \\
Nerem, Frederikse \& Hamlington (2022) & Observations (altimetry) & 1 & Very high \\
Perrette \& Mengel (2025)  & Observations (tide gauge) & 0 (spatial) & High \\
\bottomrule
\end{tabular}
\end{table}

\noindent The core conclusion of the 2026-05-10 search stands: no published
framework combines (i) direct Bayesian calibration of each component to
observational data, (ii) a structured deep-uncertainty treatment for ice
sheet instabilities, and (iii) rate-space blending with satellite-era
observations to anchor near-term projections. The 2026-09-03 update
strengthens rather than overturns that conclusion, but identifies two
papers the manuscript should engage with directly rather than simply add to
the bibliography: Darnell et al.\ (2025), whose emissions-vs-geophysical
variance decomposition is close to a rival claim (Section~4.2), and
Felikson et al.\ (2025), whose community-consensus perspective is close to
a direct statement of SPECTRUM's motivating philosophy (Section~2.2).

%% ========================================================================
\section{Action items for the manuscript}

\begin{enumerate}[leftmargin=*]
\item Cite Vermeer \& Rahmstorf (2009) --- already in \texttt{references.bib}
as \texttt{Vermeer2009} but unused --- wherever the semi-empirical class is
discussed, and consider citing Grinsted \& Christensen (2021) alongside it
if the manuscript takes an explicit position on the rate-vs-level question.
\item Add Held (2005) (and optionally Jeevanjee et al.\ 2017) as a second
anchor citation, alongside Green (2015), for the model-hierarchy/complexity
argument in the introduction.
\item Add Felikson et al.\ (2025) to the introduction or discussion as the
field-level motivation for observation-constrained projection.
\item Add an explicit sentence engaging with Darnell et al.\ (2025) where
the manuscript makes the ``emissions cuts move the median, not the tail''
claim, distinguishing SPECTRUM's fixed-horizon tail-shape claim from their
across-horizon variance decomposition.
\item Qualify the ``process models diverge from observations'' claim
(Section~1.5/1) to specify AR6 rather than IPCC assessments generally,
citing Törnqvist et al.\ (2025) as the reason for the qualification.
\item Consider citing Naughten et al.\ (2023) as the primary independent,
dynamical-modeling support for the claim that the WAIS-driven tail is
largely insensitive to emissions scenario, alongside Kopp et al.\ (2017)
and McCormack et al.\ (2026).
\item Consider a brief citation to Levermann et al.\ (2014) when introducing
the Greenland delayed-response discharge model, as the closest published
precedent for a linear, time-lagged discharge parameterization.
\end{enumerate}

%% ========================================================================
\section{Note on Felikson et al.\ (2025): why ``near-perfect conceptual match''}

Felikson et al.\ (2025) is not a competing forecasting framework in the
sense of Sections~1--2 --- it makes no projection of its own. What makes it
a near-perfect conceptual match is that it is a synthesis, by a broad and
senior cross-institutional author group spanning the leading process-based
and observational SLR communities, of exactly the argument SPECTRUM is
built to demonstrate: that large-uncertainty, physics-heavy component
projections should be narrowed using observations, that this requires
systematically quantifying uncertainty \emph{structure} (not just
magnitude) in both the observations and the models, and that doing so is
urgent ahead of the next IPCC cycle. SPECTRUM can be read as one concrete,
worked instantiation of the research program that paper calls for at the
level of an entire field. That is why it belongs in the introduction or
discussion as motivation/context, rather than in the same comparative table
as Mengel, Kopp, BRICK, FACTS, or AR6.

\end{document}
